\question{25}{
Dado un número natural $n \ge 1$, determine si es posible obtener ese número $n$ a partir de $1$ aplicando repetidamente alguna de las siguientes dos operaciones:
\begin{itemize}
  \item Duplicar y sumarle $1$.
  \item Multiplicar por $3$.
\end{itemize}

E.g.:
\begin{itemize}
  \item $n=1$: Sí (no se necesita ninguna operación).
  \item $n=3$: Sí ($1, 3$).
  \item $n=5$: No.
  \item $n=7$: Sí ($1, 3, 7$).
  \item $n=8$: No.
  \item $n=19$: Sí ($1, 3, 9, 19$).
\end{itemize}
}

\begin{enumerate}
  \item (10 décimas) Defina una relación de recurrencia que resuelva el problema. Explique en palabras qué significa la función y qué significa cada variable. Sea sucinto.
  \item (5 décimas) Dibuje y explique el grafo de dependencias. Determine cuál será la complejidad espacial de su solución.
  \item (10 décimas) Implemente una solución de programación dinámica, ya sea bottom-up o top-down, en Python, Java, C, C++, JavaScript o TypeScript. Determine la complejidad temporal de su solución.
\end{enumerate}

\bigskip
